\paragraph{2.2. Анализ адгоритмов}
\subparagraph{2.2.1}
Выразите функцию $n^3/1000 - 100n^2 - 100n + 3$ в $\Theta$-обозначениях.
$$n^3/1000 - 100n^2 - 100n + 3 = \Theta(n^3)$$

\subparagraph{2.2.2}
Рассмотрим сортировку элементов массива $A$, которая выполняется следующим образом. 
Сначала определяется наименьший элемент массива $A$, который становится на место $A[1]$. 
Затем производится поиск второго наименьшего элемента массива $A$, который становится на место элемента $A[2]$.
Этот процесс продолжается для первых $n-1$ элементов массива $A$.
Запишите псевдокод этого алгоритма, известного как сортировка выбором (selection sort).
Какой инвариант цикла сохраняется для этого алгоритма? 
Почему достаточно выполнить для первых $n-1$?
Определите время работы алгоритма в наилучшем и наихудшем случаях и запишите его в $\Theta$-обозначениях. \\

\begin{tabular}{ll}
Selection-Sort($A$) & \\
\quad for $i = 1$ to $A.length - 1$ & \\
\quad \quad $minIndex = i$ & \\ 
\quad \quad for $j = i + 1$ to $A.length$ & \\
\quad \quad \quad if $A[minIndex] > A[j]$& \\
\quad \quad \quad \quad $minIndex = j$& \\
\quad \quad $swap(A[i], A[minIndex])$ &\\
\end{tabular} \\
\\
\textbf{Инвариант цикла:}
В начале каждой итерации внешнего цикла for подмассив $A[1 \ldots i-1]$ содержит $i-1$ наименьших элементов массива $A[1 \ldots n]$, отсортированных по возрастанию. \\
\\
\textbf{Почему достаточно выполнить для первых $n-1$?}
После $n-1$ итераций внешнего цикла первые $n-1$ элементов массива отсортированы и содержат $n-1$ наименьших элементов.
Следовательно, последний элемент $A[n]$ автоматически является наибольшим, и весь массив $A[1 \ldots n]$ оказывается отсортированным.\\
\\
\textbf{Время работы алгоритма:}
Количество сравнений во внутреннем цикле не зависит от исходного порядка элементов, поэтому время работы одинаково в наилучшем и наихудшем случаях. 
Общее количество сравнений:
\[
\sum_{i=1}^{n-1} \sum_{j=i+1}^{n} 1 = \sum_{i=1}^{n-1} (n-i) = \frac{n(n-1)}{2} = \Theta(n^2).
\]
Таким образом, время работы сортировки выбором составляет $\Theta(n^2)$ в обоих случаях.

\subparagraph{2.2.3}
Вновь обратимся к алгоритму линейного поиска (см. упр. 2.1.3). 
Для скольких элементов последовательности в среднем нужно произвести проверку,
если предполагается, что все элементы массива с равной вероятностью могут иметь искомое значение?
Что происходит в наихудшем случае? Чему равно время работы алгоритма линейного поиска в среднем и в наихудшем случаях в $\Theta$-обозначениях. Обоснуйте ответ.

\textbf{Ответ:}

Предположим, что искомое значение $\upsilon$ гарантированно присутствует в массиве $A[1 \ldots n]$, и все $n$ позиций равновероятны.

\textbf{Среднее количество проверок:}
Если $\upsilon$ находится на позиции $i$, алгоритм выполняет $i$ проверок (сравнений). Математическое ожидание:
\[
\frac{1}{n} \sum_{i=1}^{n} i = \frac{1}{n} \cdot \frac{n(n+1)}{2} = \frac{n+1}{2}.
\]
Таким образом, в среднем требуется проверить $\frac{n+1}{2}$ элементов.

\textbf{Наихудший случай:}
Наихудший случай наступает, когда $\upsilon$ находится на последней позиции ($i = n$). Тогда алгоритм выполняет $n$ проверок.

\textbf{Время работы в $\Theta$-обозначениях:}
\[
\text{Средний случай: } \frac{n+1}{2} = \Theta(n), \qquad
\text{Наихудший случай: } n = \Theta(n).
\]

\subparagraph{2.2.4}
Каким образом можно модифицировать почти каждый алгоритм, чтобы получить оптимальное время работы в наилучшем случае?

\textbf{Ответ:}

Наилучший случай наступает, когда входные данные уже удовлетворяют инварианту цикла, то есть не требуют преобразований. 
Для оптимизации достаточно добавить проверку этого условия в начале цикла или внутри него, чтобы при его выполнении алгоритм завершался досрочно.